\documentclass[11pt,a4paper]{article}

\usepackage[T1]{fontenc}
\usepackage[utf8]{inputenc}
\usepackage[catalan]{babel}
\usepackage{lmodern}
\usepackage[a4paper,margin=2.5cm]{geometry}
\usepackage{graphicx}
\usepackage{xcolor}
\usepackage{booktabs}
\usepackage{tabularx}
\usepackage{longtable}
\usepackage{enumitem}
\usepackage{amssymb}
\usepackage{listings}
\usepackage[unicode]{hyperref}

\hypersetup{
  colorlinks=true,
  linkcolor=black!70!blue,
  urlcolor=blue!60!black,
  pdftitle={Informe del projecte: El torneig dels mags},
  pdfauthor={Gerard Casacuberta, Diego Suarez, Mikolaj Jaworski}
}

\lstset{
  basicstyle=\ttfamily\small,
  breaklines=true,
  frame=single,
  columns=fullflexible,
  keepspaces=true
}

\setlist{itemsep=2pt,topsep=4pt}

\title{\textbf{El torneig dels mags}\\[0.4em]
\large Informe del projecte de Sistemes Multijugador\\[0.3em]
\normalsize Mòdul de Gestió d'Usuaris i Mòdul d'Interacció}
\author{Gerard Casacuberta Albert \and Juan Diego Suárez P. \and Mikolaj Jaworski}
\date{\today}

\begin{document}

\maketitle

\begin{abstract}
\noindent
Aquest informe documenta les dues parts del projecte i com es relacionen entre elles.
El \textbf{Mòdul de Gestió d'Usuaris} és una aplicació Node.js amb plantilles Eta que
permet registrar-se, iniciar i tancar sessió, editar el perfil, canviar la contrasenya,
recuperar-la per correu i eliminar el compte. El \textbf{Mòdul d'Interacció} és el joc
multijugador: un taulell isomètric en PixiJS\ amb personatges de sis cases que competeixen
per aconseguir la copa del laberint, sincronitzat en temps real amb Colyseus.
Els dos mòduls comparteixen el mateix procés de servidor i el mateix testimoni JWT, de
manera que la identitat creada al mòdul d'usuaris és la que autoritza l'entrada a la partida.
\end{abstract}

\tableofcontents
\newpage

% ---------------------------------------------------------------------------
\section{Objectiu i abast}

El projecte consisteix en un joc multijugador en xarxa i en el sistema d'usuaris que el
sustenta. L'objectiu del joc és recórrer un laberint isomètric de $60 \times 60$ caselles,
recollir boles que amplien el moviment i arribar abans que ningú a la copa situada a
l'extrem oposat del mapa.

El treball s'ha dividit en dos mòduls que es poden desplegar per separat però que estan
dissenyats per funcionar junts:

\begin{description}[leftmargin=1.2cm,style=nextline]
  \item[Mòdul de Gestió d'Usuaris] Node.js, Express, plantilles Eta i SQLite.
        És la porta d'entrada: cap jugador no pot connectar-se a una partida sense un
        compte vàlid i sense sessió activa.
  \item[Mòdul d'Interacció] Client PixiJS\ i servidor de sales Colyseus.
        És la partida en si: lobby, selecció d'equip, moviment, disputes per casella,
        millores, condició de victòria i reconnexió a una altra partida.
\end{description}

A més de descriure cada mòdul, l'informe dedica una atenció especial a la part
multijugador de cadascun: la identitat compartida i la moderació al primer, i la
sincronització d'estat, la latència i l'arbitratge de clics al segon.

% ---------------------------------------------------------------------------
\section{Arquitectura general}

El repositori conté dos projectes npm independents, sense \texttt{package.json} arrel:

\begin{verbatim}
  navegador
     |
     |  HTTPS (REST + pàgines Eta)        WSS (Colyseus)
     v                                        v
  +-----------------------------------------------------+
  |  servidor Node.js  (Express + Colyseus, port 2567)  |
  |  index.js  ->  rutes, plantilles, API REST          |
  |  GameRoom.js -> sales, estat, moviment, partida     |
  |  db/  -> SQLite (better-sqlite3)                    |
  +-----------------------------------------------------+
     ^
     |  fitxers estàtics + build Vite
  +-----------------------------------------------------+
  |  client Vite/PixiJS  (dev 5173, nginx 8080)         |
  +-----------------------------------------------------+
\end{verbatim}

El servidor allotja en un sol procés HTTP i WebSocket: Express atén les pàgines i l'API
REST, i el transport WebSocket de Colyseus atén les sales de joc. El client obté l'adreça
del servidor de \texttt{/config.js} (variable \texttt{window.\_\_SERVER\_URL\_\_}, generada
pel contenidor); en desenvolupament fa servir \texttt{http://<host>:2567}.

Les decisions estructurals més rellevants són:

\begin{itemize}
  \item \textbf{Moviment autoritatiu al servidor.} El client només envia el destí; el
        servidor valida, calcula el camí i retransmet la ruta. El client no pot teleportar-se.
  \item \textbf{Un únic origen de veritat per a l'usuari.} La taula \texttt{users} de
        SQLite alimenta l'autenticació web i l'entrada a les sales; no hi ha dues bases de dades.
  \item \textbf{Testimoni JWT com a contracte entre mòduls.} El mateix secret signa i
        verifica el token a l'API HTTP i a \texttt{GameRoom.onAuth}.
  \item \textbf{Sense ORM ni framework d'UI.} SQL parametritzat amb \texttt{better-sqlite3}
        i plantilles Eta, per mantenir el codi petit i auditable.
\end{itemize}

% ---------------------------------------------------------------------------
\section{Mòdul de Gestió d'Usuaris}

\subsection{Tecnologies i estructura}

Node.js $\geq 20.11$, Express 4, Eta 4 com a motor de plantilles, \texttt{better-sqlite3}
com a accés directe a SQLite, \texttt{argon2} per als hash de contrasenya,
\texttt{jsonwebtoken} per a la sessió, \texttt{multer} per a l'avatar i l'API REST de Brevo
per al correu de recuperació.

\begin{itemize}[nosep]
  \item \texttt{server/index.js}: rutes de pàgines i de l'API, middleware global.
  \item \texttt{server/GameRoom.js}: sala de joc (pertany al mòdul d'interacció, però
        comparteix procés i autenticació).
  \item \texttt{server/utils.js}: JWT, \texttt{attachUser}, \texttt{auth}, validacions comunes.
  \item \texttt{server/db/}: connexió, esquema i consultes parametritzades.
  \item \texttt{server/passwordPolicy.js} i \texttt{passwordRecovery.js}: política de
        contrasenyes i flux de recuperació.
  \item \texttt{server/templates/*.eta}: \texttt{home}, \texttt{login}, \texttt{register},
        \texttt{profile}, \texttt{delete}, \texttt{banned}, \texttt{404}.
\end{itemize}

\subsection{Model de dades}

\begin{lstlisting}[language=SQL]
CREATE TABLE users (
    id INTEGER PRIMARY KEY,
    username VARCHAR(63) UNIQUE,
    hash VARCHAR(255),           -- Argon2id, mai la contrasenya
    email VARCHAR(255) UNIQUE,
    birthdate TEXT,
    nickname VARCHAR(150) DEFAULT '',
    avatar VARCHAR(255) DEFAULT '/icons/default.svg',
    isBanned BOOLEAN DEFAULT 0,
    reason TEXT
);

CREATE TABLE password_resets (
    user_id INTEGER PRIMARY KEY REFERENCES users(id) ON DELETE CASCADE,
    token_hash TEXT NOT NULL UNIQUE,
    expires_at INTEGER NOT NULL
);
\end{lstlisting}

L'esquema es crea de forma idempotent a l'arrencada (\texttt{db/create\_db.sql}), de manera
que \texttt{npm start}, \texttt{npm run dev} i \texttt{npm run db} funcionen sobre una base
buida. La ruta es pot sobreescriure amb \texttt{DB\_PATH}.

\subsection{Registre}

\texttt{POST /api/register} valida, en aquest ordre: camps obligatoris (usuari, correu, data
de naixement i contrasenya), longitud mínima de l'usuari (3 caràcters, com al mòdul PHP
anterior), format del correu, data de naixement real i edat mínima de 16 anys, coincidència
de la confirmació si s'envia i política de contrasenya. La política és compartida amb la
recuperació i el canvi de contrasenya: mínim 12 caràcters Unicode, sense exigir majúscules
ni símbols, i rebuig de contrasenyes filtrades consultant Pwned Passwords amb
\textit{k-anonymity} (només s'envien els cinc primers caràcters del SHA-1). Si el servei
extern no respon, es permet continuar i només es registra el codi \texttt{HIBP\_UNAVAILABLE}.

La contrasenya s'emmagatzema amb Argon2id i la inserció aprofita les restriccions
\texttt{UNIQUE} de SQLite per respondre amb un 409 quan l'usuari o el correu ja existeixen.
El registre retorna el token i deixa la cookie de sessió preparada.

\subsection{Sessió i autenticació}

El token és un JWT de dues hores amb \texttt{id} i \texttt{username}. Viatja en una cookie
\texttt{httpOnly}, \texttt{sameSite=lax} i \texttt{secure} quan
\texttt{COOKIE\_SECURE=production}, i també es retorna al cos de la resposta perquè el
client Vite el pugui reutilitzar en un altre origen. El middleware \texttt{attachUser}
verifica la cookie a cada petició; \texttt{auth} protegeix les pàgines i les rutes
d'API (retornant \texttt{401} JSON a \texttt{/api/*} i una redirecció a \texttt{/login} a
les pàgines).

\subsection{Perfil i avatar}

\texttt{GET /profile} renderitza la plantilla \texttt{profile.eta} amb les dades de
l'usuari. \texttt{PUT /api/profile} accepta \texttt{multipart/form-data} amb un sobrenom i,
opcionalment, una imatge: límit de 2 MB, tipus permesos JPEG, PNG i WebP, nom aleatori
generat amb \texttt{randomUUID()} (no s'usa mai el nom original) i esborrat de l'avatar
anterior si era una pujada nostra. Els fitxers viuen a \texttt{UPLOAD\_DIR} i es publiquen
a \texttt{/uploads}.

\subsection{Canvi de contrasenya}

\texttt{POST /api/password} demana la contrasenya actual, la nova i la confirmació. Comprova
els camps, la coincidència, la política de 12 caràcters i, finalment, verifica la
contrasenya actual amb Argon2 abans d'actualitzar el hash. És el darrer mòdul heretat del
panell PHP i comparteix el mateix validador que el registre i la recuperació.

\subsection{Recuperació de contrasenya}

El flux evita revelar si un correu existeix: la resposta és sempre la mateixa i es retorna
abans de consultar la base de dades o de contactar amb Brevo. Els límits són 5 peticions
per IP cada 15 minuts i, per correu, 1 per minut i 3 per hora (la clau és el SHA-256 del
correu, no el correu en clar). Es genera un token de 32 bytes aleatoris, es guarda només el
seu SHA-256 amb una caducitat de 30 minuts i s'envia l'enllaç per l'API REST de Brevo.
L'enllaç apunta al client (\texttt{FRONTEND\_URL\#reset-password=<token>}); el servidor
només redirigeix \texttt{/forgot} cap allà. El consum del token és una transacció
\texttt{immediate} de SQLite: es revalida, s'actualitza el hash i s'esborren els tokens de
l'usuari, de manera que dues peticions simultànies no poden canviar la contrasenya dues
vegades. Els errors de correu es registren només amb codis
(\texttt{BREVO\_HTTP\_4xx}, \texttt{BREVO\_TIMEOUT}, etc.).

\subsection{Eliminació de compte}

\texttt{POST /api/delete} exigeix la contrasenya actual, elimina en una transacció els
tokens de recuperació i l'usuari, neteja la cookie, esborra l'avatar pujat i respon
\texttt{204}. Les rutes de pàgina \texttt{/delete} i \texttt{/profile} comproven que
l'usuari existeixi a la base de dades, no només que el token sigui vàlid.

\subsection{Baneig i moderació}

La columna \texttt{isBanned} i el motiu associat s'apliquen en un middleware global: un
compte bloquejat només pot veure \texttt{/banned} (i tancar sessió via
\texttt{/api/logout}); qualsevol altra pàgina el redirigeix i qualsevol altra crida a
l'API retorna \texttt{403} amb el motiu i la ruta de destí. Això evita l'estat sense
sortida: si el baneig bloquegés també el tancament de sessió, l'usuari quedaria atrapat.

\subsection{La part multijugador d'aquest mòdul}

Encara que el mòdul no és cap joc, és el que fa possible que diverses persones juguin
alhora:

\begin{itemize}
  \item \textbf{Identitat única per a tots els canals.} El mateix JWT autoritza les
        peticions HTTP i l'entrada a la sala WebSocket, encara que el client i el servidor
        estiguin en dominis diferents (\texttt{CLIENT\_ORIGIN} i CORS amb credencials).
  \item \textbf{Una sola llista de jugadors.} El registre, el perfil i el baneig són els
        mateixos per a tothom; no existeix un ``usuari del joc'' separat de l'usuari web.
  \item \textbf{Concurrència.} \texttt{better-sqlite3} és síncron i SQLite té un sol
        escriptor; juntament amb les transaccions, això evita condicions de cursa en
        registres duplicats i en el consum de tokens. Els límits de peticions són en
        memòria del procés, amb neteja mandrosa perquè la memòria no creixi sense control.
  \item \textbf{Privacitat entre jugadors.} La recuperació de contrasenya no confirma
        l'existència d'un compte, i les rutes d'API no retornen mai el hash ni el correu
        d'altres usuaris.
  \item \textbf{Moderació efectiva.} El baneig s'aplica a les pàgines, a l'API i al
        \texttt{onAuth} de la sala, de manera que un compte bloquejat no pot continuar
        jugant amb una pestanya oberta.
\end{itemize}

% ---------------------------------------------------------------------------
\section{Mòdul d'Interacció}

\subsection{Tecnologies i pantalles}

El client és una aplicació Vite amb PixiJS\ v8 per al renderitzat, \texttt{pixi-tiledmap}
per carregar el mapa de Tiled i \texttt{colyseus.js} per a la xarxa. Les pantalles són el
menú (entrar o registrar-se, amb l'opció d'indicar un identificador de sala), la
recuperació de contrasenya, el lobby d'equips, la partida i el resultat.

\subsection{Mapa i renderitzat isomètric}

El mapa viu a \texttt{client/public/assets/map.tmj}: $60 \times 60$ caselles de
$232 \times 110$ píxels amb les capes \texttt{ground}, \texttt{grass}, \texttt{sanctuary},
\texttt{walls}, \texttt{details} i \texttt{collision}. Les propietats del mapa guarden les
posicions de sortida dels sis equips, la copa a $[55,55]$ i vuit boles de millora.

El servidor llegeix el mateix fitxer en arrencar: en deriva la mida de la graella i la
llista de caselles bloquejades de la capa \texttt{collision}, o sigui que el mapa és
l'única font de veritat per al client i per al servidor. Per això cal reiniciar el servidor
després de modificar-lo.

Al client, \texttt{renderTerrainSurfaces()} recorre les capes de terra i crea un sol
\texttt{Sprite} per casella, mantenint l'ordre de pintat entre herba, jardí i santuari; si
cada capa es dibuixés per separat, les cares laterals d'un cub taparien el terra veí i el
mapa semblaria elevat. Els personatges són fulls d'animació de sis cases
(Gryffindor, Ravenclaw, Slytherin, Hufflepuff, Beauxbatons i Durmstrang). Cada full té tres
files (perfil, esquerra/dreta, davant/darrere) i un nombre de fotogrames diferent; el codi
retalla cada fotograma pels píxels opacs i l'escala a $88 \times 128$ perquè tots els
personatges tinguin la mateixa mida en pantalla. La direcció del moviment tria la fila, i
la càmera escala el món per cobrir la finestra i segueix el jugador local.

\subsection{Lobby, equips i inici de partida}

En entrar, el client crea una sala nova o s'uneix a una d'existent amb
\texttt{joinById}. El servidor admet fins a 6 jugadors; cadascun tria una casa i dues
persones no poden triar la mateixa. La partida comença quan tots els jugadors han triat
equip i se supera el mínim (1 en desenvolupament, 2 en producció). En aquell moment la sala
s'estableix com a \texttt{locked}: no s'hi pot entrar a mig joc, cosa que garanteix que tots
els clients vegin el mateix inici.

\subsection{Moviment autoritatiu i sincronització}

El protocol és intencionadament curt:

\begin{enumerate}[nosep]
  \item El jugador fa clic en una casella. El client comprova que sigui dins de l'abast,
        que no estigui bloquejada ni ocupada, i envia
        \texttt{moveTo \{x, y, clickedAt: serverNow()\}}.
  \item El servidor comprova la fase, que el jugador no estigui caminant ja, els límits, la
        col·lisió, l'ocupació i el resultat de la cerca en amplada (BFS) sobre la graella.
        La ruta ha de tenir almenys dues caselles i no superar \texttt{min(moveRange, 7)}.
  \item Si és vàlida, retransmet a tothom:
\end{enumerate}

\begin{lstlisting}
{ id, path: [{x, y}, ...], startTime: Date.now(), msPerTile: 300 }
\end{lstlisting}

Cada client interpola el personatge amb \texttt{serverNow() = Date.now() + offset}, on
\texttt{offset} és la mitjana de les vuit darreres mostres de \texttt{ping}/\texttt{pong}
(la marca de temps del servidor menys la del client, ajustada pel temps d'anada i tornada).
Així tothom veu la mateixa animació encara que els rellotges estiguin desfasats. El client
no fa predicció local: la posició sempre arriba del servidor. A canvi de més latència
percebuda, s'elimina qualsevol possibilitat de desincronització o de trampa de posició.

\subsection{Disputes per una casella}

Dos jugadors poden clicar la mateixa casella gairebé alhora. El servidor no decideix amb el
primer missatge que arriba, perquè la latència afavoriria uns jugadors i no altres:

\begin{itemize}
  \item Cada destí té una \emph{reclamació} amb una finestra d'arbitratge de 250 ms,
        ampliada fins a la meitat del RTT més alt dels clients (màxim 1500 ms).
  \item Els candidats s'ordenen pel \texttt{clickedAt} estimat i, si empaten, per l'ordre
        d'arribada. El servidor rebutja marques de temps massa antigues o futures.
  \item El guanyador recorre la ruta sencera. Els perdedors avancen fins a la casella
        anterior a la disputada; en arribar-hi es retransmet \texttt{duel} i retrocedeixen
        un 30\% dels passos (mínim un) amb una animació de retrocés.
  \item Els temporitzadors de recollida de boles i de la copa s'associen a cada pas, de
        manera que només el qui trepitja la casella rep la millora o la victòria.
\end{itemize}

Aquest disseny converteix la latència en un problema del servidor i no del jugador: el
resultat és el mateix per a tothom, i els clients només reprodueixen el que s'ha decidit.

\subsection{Millores i condició de victòria}

Hi ha vuit boles al mapa. Recollir-ne una augmenta l'abast de moviment en una casella, de 3
a un màxim de 7, i desapareix de l'estat compartit; a partir del màxim, queden disponibles
per als altres. Guanya qui arriba primer a la casella de la copa. En aquell moment el
servidor atura les disputes, cancel·la els temporitzadors pendents, fixa el guanyador i
retransmet l'estat \texttt{finished}; tots els clients veuen el mateix desenllaç.

\subsection{HUD, ping i resultat}

La interfície mostra l'abast actual, el panell de ping de cada jugador (canvia de color
segons la latència, es pot obrir amb el Tab) i la pantalla de resultat amb el nom del
guanyador. Des d'allà es pot tornar a jugar (nova sala) o tancar la sessió. En sortir d'una
sala, el client es torna a unir automàticament si l'usuari tria tornar a jugar.

\subsection{La part multijugador del joc}

\begin{itemize}
  \item \textbf{Estat compartit declaratiu.} Colyseus\ sincronitza un \texttt{State} amb la
        fase, la graella de col·lisions, el mapa de jugadors, les boles, la copa i el
        guanyador. El client només dibuixa el que rep; no hi ha estat duplicat a cada costat.
  \item \textbf{Reconciliació temporal.} Les rutes porten \texttt{startTime} i
        \texttt{msPerTile}, i les animacions de duel porten la seva pròpia marca i durada.
        El rellotge estimat permet que una animació de 300 ms per casella es vegi igual a
        tots els navegadors.
  \item \textbf{Justícia.} L'arbitratge per marca de clic, la finestra adaptada al RTT i el
        retrocés parcial dels perdedors eviten que el guanyador sigui simplement el que té
        menys latència.
  \item \textbf{Integritat.} El servidor valida cada moviment (fase, límits, col·lisions,
        ocupació, ruta i abast) i retransmet només rutes legals. Un client manipulat no pot
        saltar-se el mapa ni moure's més caselles de les permeses.
  \item \textbf{Presència i reconnexió.} En sortir, la sala allibera l'equip i cancel·la les
        disputes obertes del jugador. Si la connexió es perd, el client recarrega la pàgina.
  \item \textbf{Lobby sense sales públiques.} L'identificador de sala es pot compartir per
        pantalla; no hi ha navegador de sales, cosa que simplifica el disseny i evita
        llistats costosos.
\end{itemize}

% ---------------------------------------------------------------------------
\section{Com es connecten els dos mòduls}

La connexió és el testimoni JWT i un únic procés de servidor:

\begin{enumerate}
  \item L'usuari s'autentica contra el mòdul de gestió (\texttt{/api/login} o
        \texttt{/api/register}). El servidor signa un token i el desa en una cookie
        \texttt{httpOnly}; el client també el rep al cos de la resposta.
  \item El client del joc passa el token com a opció de connexió a Colyseus:
        \texttt{create("game", \{token\})} o \texttt{joinById(id, \{token\})}.
  \item \texttt{GameRoom.onAuth} verifica el token amb el mateix \texttt{JWT\_SECRET},
        rebutja els comptes bloquejats i impedeix que el mateix usuari entri dues vegades.
        Si la verificació falla, la connexió es tanca abans de crear cap jugador.
  \item A partir d'aquí, el nom d'usuari acompanya el jugador a l'estat de la sala i és el
        que es mostra sobre el personatge.
\end{enumerate}

També hi ha dos llaços de tornada del joc cap al mòdul d'usuaris:

\begin{itemize}
  \item El baneig es decideix a la base de dades compartida i s'aplica tant a les rutes
        HTTP com a l'entrada a la sala.
  \item L'enllaç de recuperació de contrasenya apunta al client Vite
        (\texttt{FRONTEND\_URL}), perquè el formulari de nova contrasenya viu a la mateixa
        interfície que el joc; el servidor només redirigeix \texttt{/forgot}.
\end{itemize}

El desplegament manté aquesta separació: \texttt{docker compose} aixeca el servidor al port
2567 i el client (nginx) al 8080; a Render hi ha dos serveis Docker, un disc persistent d'1
GB per a la base de dades i els avatars al servidor, i el client genera
\texttt{/config.js} a partir de \texttt{SERVER\_URL}. El servidor de joc només necessita el
\texttt{map.tmj} del client, que es copia a la imatge.

% ---------------------------------------------------------------------------
\section{Proves i qualitat}

El projecte no té un \texttt{npm test} al client, però el servidor inclou 35 proves amb
\texttt{node --test}:

\begin{itemize}
  \item \textbf{Política de contrasenyes}: longitud Unicode, mínim de 12 caràcters,
        contrasenyes filtrades, k-anonymity i tolerància a caigudes de Pwned Passwords.
  \item \textbf{Recuperació}: emmagatzematge només del hash, caducitat, un sol ús, consum
        concurrent, límits per IP i per correu, mateixa resposta per a correus inexistents i
        errors de Brevo sense dades sensibles.
  \item \textbf{Mapa}: grandària, accessibilitat de totes les caselles, absència de zones
        obertes de $3 \times 3$, sortides $2 \times 3$ i sincronia entre els fitxers de Tiled.
  \item \textbf{Partida}: reinici d'estat entre rondes, victòria i sortida de clients, i una
        prova de xarxa real amb dos clients Colyseus\ que veuen el mateix guanyador i tornen
        a jugar.
\end{itemize}

A més, la compilació del client (\texttt{npm run build}) i el desplegament amb Docker
serveixen de comprovació d'integració, i les rutes principals s'han provat manualment amb
comptes reals (registre, sessió, canvi de contrasenya, baneig i partida de dos jugadors).

% ---------------------------------------------------------------------------
\section{Distribució de tasques}

El treball s'ha organitzat per funcionalitats verticals i diverses tasques s'han fet en
parelles o en trucades de Discord, tal com reflecteix la taula. La distribució és
aproximadament equilibrada: cadascú lidera un bloc i participa en la integració i les
proves.

\begin{longtable}{@{}p{7.6cm}ccc@{}}
\toprule
\textbf{Tasca} & \textbf{Gerard} & \textbf{Diego} & \textbf{Mikolaj} \\
\midrule
\endhead
Disseny funcional, model de dades i esquema SQL & \checkmark & \checkmark & \checkmark \\
Backend d'autenticació (registre, login, logout) & \checkmark & & \checkmark \\
Plantilles Eta, estils i accessibilitat de les pàgines & & \checkmark & \checkmark \\
Política de contrasenyes i HIBP & & & \checkmark \\
Recuperació de contrasenya i integració amb Brevo & \checkmark & & \checkmark \\
Perfil, pujada d'avatar i neteja de fitxers & & \checkmark & \\
Canvi de contrasenya & \checkmark & & \\
Eliminació de compte i transaccions & \checkmark & \checkmark & \\
Baneig, motius i pantalla de compte bloquejat & & \checkmark & \checkmark \\
Mapa de Tiled, capes i col·lisions & \checkmark & & \checkmark \\
Moviment autoritatiu, BFS i validació al servidor & \checkmark & & \\
Renderitzat isomètric, càmera i ordre de pintat & & \checkmark & \\
Personatges, fulls d'animació i direccions & & & \checkmark \\
Lobby, selecció d'equip i inici de partida & \checkmark & \checkmark & \\
Sincronització de rellotge, ping i interpolació & \checkmark & \checkmark & \\
Disputes per casella, arbitratge i retrocés & \checkmark & & \checkmark \\
Boles de millora, abast i condició de victòria & & \checkmark & \checkmark \\
HUD, panell de ping i pantalla de resultat & & & \checkmark \\
Integració JWT entre HTTP i WebSocket & \checkmark & \checkmark & \\
Docker, Render i variables d'entorn & & \checkmark & \checkmark \\
Proves automatitzades i comprovacions manuals & \checkmark & \checkmark & \checkmark \\
Documentació i informe & \checkmark & \checkmark & \checkmark \\
\bottomrule
\end{longtable}

% ---------------------------------------------------------------------------
\section{Enllaços}

Els enllaços definitius s'afegiran en lliurar el projecte; de moment es conserven com a
places de referència:

\begin{enumerate}
  \item \textbf{Codi font (ambdós mòduls)}:
        \url{https://example.com/torneig-mags-codi.zip}
  \item \textbf{Joc en línia (client)}:
        \url{https://example.com/torneig-mags}
  \item \textbf{Aplicació en línia (gestió d'usuaris)}:
        \url{https://example.com/torneig-mags-usuaris}
\end{enumerate}

\noindent\emph{Nota: substituïu aquestes adreces per les definitives abans de l'entrega.}

% ---------------------------------------------------------------------------
\section{Limitacions i treball futur}

\begin{itemize}
  \item \textbf{Sense predicció local.} El moviment és totalment autoritatiu, cosa que
        garanteix la coherència però afegeix la meitat del RTT a cada clic. Una predicció
        amb reconciliació seria la millora natural.
  \item \textbf{Tokens no revocables.} El JWT és sense estat i dura dues hores. El baneig
        immediat es comprova a la base de dades a les rutes HTTP; a la sala es valida el
        token en el moment de la connexió, no a cada moviment. Una llista de revocació o
        tokens més curts ho resoldrien del tot.
  \item \textbf{Límits de peticions en memòria.} Funcionen bé amb un sol procés; amb
        diverses instàncies caldria un magatzem compartit.
  \item \textbf{Sense historial de partides.} Ni rànquing ni estadístiques persistents;
        l'estat de la sala viu només mentre duri la partida.
  \item \textbf{Una sola sala per codi.} L'identificador és l'habitació de Colyseus; no hi
        ha llistat de sales públiques ni aparellament automàtic.
  \item \textbf{Avatar i sobrenom només a la web.} El joc mostra el nom d'usuari; portar el
        sobrenom i l'avatar a l'estat compartit seria un canvi petit i visible.
\end{itemize}

\section{Conclusions}

El resultat és un joc multijugador complet amb un sistema d'usuaris real: la mateixa
identitat serveix per entrar a la web i a la partida, el servidor decideix tot el que
afecta la partida i tots els clients veuen el mateix món encara que tinguin latències
diferents. La separació en dos mòduls ha permès treballar en paral·lel (interfície i
autenticació d'una banda, mapa i sincronització de l'altra) sense duplicar dades ni
lògica, i ha deixat el camí obert a millores com la predicció de moviment o l'historial
de partides.

\end{document}
